% =============================================================
% TEMPLATE BÁO CÁO MÔN HỌC - HCMUS
% Chỉ cần sửa phần THÔNG TIN BÁO CÁO bên dưới, sau đó viết nội
% dung trong thư mục content/. Danh sách sinh viên, giảng viên
% sửa trong config/frontpage.tex.
% Biên dịch: nút xanh trong VS Code, hoặc latexmk -pdf main.tex
% =============================================================

% ---------------- THÔNG TIN BÁO CÁO ----------------
\newcommand{\subjectname}{Tên môn học}            % hiện ở header và trang bìa
\newcommand{\subjectcode}{CSC00000}               % mã môn
\newcommand{\reportname}{BÁO CÁO ĐỒ ÁN}          % loại báo cáo: đồ án, bài tập lớn, thực hành...
\newcommand{\reporttitle}{Chủ đề: Tên đề tài}
\newcommand{\classname}{Tên lớp}
\newcommand{\semester}{Học kỳ 1, năm học 2025--2026}
\newcommand{\reportdate}{TP Hồ Chí Minh, tháng 1 năm 2026}

% ---------------- THÔNG TIN TRƯỜNG ----------------
\newcommand{\head}{Đại học Quốc gia Thành phố Hồ Chí Minh}
\newcommand{\university}{Trường Đại học Khoa học tự nhiên}
\newcommand{\faculty}{Khoa Công nghệ thông tin}

% =============================================================
% PREAMBLE: gói lệnh, header/footer, màu sắc, style code.
% Thông tin riêng của từng báo cáo (môn học, tên đề tài, ...)
% nằm ở main.tex, KHÔNG sửa trong file này.
% =============================================================
\documentclass[12pt]{article}

% ---------------- Ngôn ngữ, khổ giấy, giãn dòng ----------------
\linespread{1.3}   % giãn dòng toàn văn bản (1 = dòng đơn)
\usepackage[utf8]{inputenc}
\usepackage[vietnamese]{babel}
\usepackage[a4paper,top=1in,bottom=1in,left=1.25in,right=1.25in,marginparwidth=2cm]{geometry}
\usepackage{indentfirst}

% ---------------- Toán, bảng, hình, danh sách ----------------
\usepackage{amsmath}
\usepackage{amsfonts}
\usepackage{array}
\usepackage{tabularx}
\usepackage{graphicx}
\usepackage{float}
\usepackage{subcaption}
\usepackage[inline, shortlabels]{enumitem}
\usepackage{xcolor}
\usepackage{tikz}
\newcolumntype{C}{>{\centering\arraybackslash}X}

% ---------------- Liên kết trong PDF ----------------
\usepackage[colorlinks = true,
            linkcolor = black,
            urlcolor  = orange!30!black,
            citecolor = black,
            anchorcolor = black]{hyperref}
\hypersetup{pdftitle={\subjectname}}
\addto\captionsvietnamese{\renewcommand{\refname}{Tài liệu tham khảo}}

% ---------------- Header / footer ----------------
\usepackage{lastpage}
\usepackage{fancyhdr}
\pagestyle{fancy}
\setlength{\headheight}{30pt}
\addtolength{\topmargin}{-4pt}
\fancyhead{}
\fancyhead[L]{
    \begin{tabular}{rl}
        \begin{picture}(25,15)(0,0)
    \put(0,-8){\includegraphics[scale = 0.025]{images/hcmus-logo.png}}
       \end{picture}&\hspace{-.5cm}
        \raggedleft\begin{tabular}{l}
            \university\\
            \faculty
        \end{tabular}
    \end{tabular}
}
\fancyhead[R]{\subjectname}
\fancyfoot{}
\fancyfoot[R]{Trang {\thepage}/\pageref{LastPage}}

% ---------------- Màu cho code ----------------
\definecolor{parenthesesblue}{rgb}{0.251,0.4706,0.949}
\definecolor{codepurple}{rgb}{0.652,0.148,0.6432}
\definecolor{stringgreen}{rgb}{0.3137,0.6314,0.3098}
\definecolor{variableorange}{rgb}{0.597,0.4087,0.003}
\definecolor{delimorange}{rgb}{0.7562,0.5176,0.0038}
\definecolor{commentgray}{rgb}{0.75,0.75,0.75}
\definecolor{backcolour}{rgb}{0.97,0.97,0.97}

% ---------------- Code (listings), giả mã (algorithm2e) ----------------
\usepackage{listings}
\usepackage[most]{tcolorbox}
\newtcolorbox{pseudocodebox}{
    colback=yellow!10,
    fonttitle=\bfseries\scshape,
    boxrule=0pt,frame hidden,sharp corners,enhanced
}
\usepackage[ruled,vlined,noline]{algorithm2e}
\makeatletter
\def\addToLiterate#1{\edef\lst@literate{\unexpanded\expandafter{\lst@literate}\unexpanded{#1}}}
\lst@Key{add to literate}{}{\addToLiterate{#1}}
\makeatother
\lstdefinelanguage{pseudocode}{
  morekeywords={function,if,then,else,while,return,for,do,end},
  sensitive=false,
  morecomment=[l]{//},
  morestring=[b]"
}
\lstdefinestyle{mystyle}{
  backgroundcolor=\color{backcolour},
  commentstyle=\it\color{commentgray},
  keywordstyle=\color{codepurple},
  stringstyle=\color{stringgreen},
  basicstyle=\ttfamily\footnotesize,
  captionpos=b,
  keepspaces=true,
  showspaces=false,
  showstringspaces=false,
  showtabs=false,
  tabsize=2
}
\lstset{
    language = Python,
    style=mystyle,
    literate=%
	% Vần a
		{á}{{\'a}}1
		{à}{{\`a}}1
		{ạ}{{\d a}}1
		{ả}{{\h a}}1
		{ã}{{\~ a}}1
		%
		{Á}{{\'A}}1
		{À}{{\`A}}1
		{Ạ}{{\d A}}1
		{Ả}{{\h A}}1
		{Ã}{{\~ A}}1
	%
	% Vần ă
		{ă}{{\u a}}1
		{ắ}{{\'\abreve }}1
		{ằ}{{\`\abreve }}1
		{ặ}{{\d \abreve }}1
		{ẳ}{{\h \abreve }}1
		{ẵ}{{\~\abreve }}1
		%
		{Ă}{{\u A}}1
		{Ắ}{{\'\ABREVE }}1
		{Ằ}{{\`\ABREVE }}1
		{Ặ}{{\d \ABREVE }}1
		{Ẳ}{{\h \ABREVE }}1
		{Ẵ}{{\~\ABREVE }}1
	%
	% Vần â
		{â}{{\^ a}}1
		{ấ}{{\'\acircumflex }}1
		{ầ}{{\`\acircumflex }}1
		{ậ}{{\d \acircumflex }}1
		{ẩ}{{\h \acircumflex }}1
		{ẫ}{{\~\acircumflex }}1
		%
		{Â}{{\^ A}}1
		{Ấ}{{\'\ACIRCUMFLEX }}1
		{Ầ}{{\`\ACIRCUMFLEX }}1
		{Ậ}{{\d \ACIRCUMFLEX }}1
		{Ẩ}{{\h \ACIRCUMFLEX }}1
		{Ẫ}{{\~\ACIRCUMFLEX }}1
	%
	% Vần đ
		{đ}{{\dj }}1
		{Đ}{{\DJ }}1
	%
	% Vần e
		{é}{{\'e}}1
		{è}{{\`e}}1
		{ẹ}{{\d e}}1
		{ẻ}{{\h e}}1
		{ẽ}{{\~ e}}1
		%
		{É}{{\'E}}1
		{È}{{\`E}}1
		{Ẹ}{{\d E}}1
		{Ẻ}{{\h E}}1
		{Ẽ}{{\~ E}}1
	%
	% Vần ê
		{ê}{{\^e}}1
		{ế}{{\'\ecircumflex }}1
		{ề}{{\`\ecircumflex }}1
		{ệ}{{\d \ecircumflex }}1
		{ể}{{\h \ecircumflex }}1
		{ễ}{{\~\ecircumflex }}1
		%
		{Ê}{{\^E}}1
		{Ế}{{\'\ECIRCUMFLEX }}1
		{Ề}{{\`\ECIRCUMFLEX }}1
		{Ệ}{{\d \ECIRCUMFLEX }}1
		{Ể}{{\h \ECIRCUMFLEX }}1
		{Ễ}{{\~\ECIRCUMFLEX }}1
	%
	% Vần i
		{í}{{\'i}}1
		{ì}{{\`\i }}1
		{ị}{{\d i}}1
		{ỉ}{{\h i}}1
		{ĩ}{{\~\i }}1
		%
		{Í}{{\'I}}1
		{Ì}{{\`I}}1
		{Ị}{{\d I}}1
		{Ỉ}{{\h I}}1
		{Ĩ}{{\~I}}1
	%
	% Vần o
		{ó}{{\'o}}1
		{ò}{{\`o}}1
		{ọ}{{\d o}}1
		{ỏ}{{\h o}}1
		{õ}{{\~o}}1
		%
		{Ó}{{\'O}}1
		{Ò}{{\`O}}1
		{Ọ}{{\d O}}1
		{Ỏ}{{\h O}}1
		{Õ}{{\~O}}1
	%
	% Vần ô
		{ô}{{\^o}}1
		{ố}{{\'\ocircumflex }}1
		{ồ}{{\`\ocircumflex }}1
		{ộ}{{\d \ocircumflex }}1
		{ổ}{{\h \ocircumflex }}1
		{ỗ}{{\~\ocircumflex }}1
		%
		{Ô}{{\^O}}1
		{Ố}{{\'\OCIRCUMFLEX }}1
		{Ồ}{{\`\OCIRCUMFLEX }}1
		{Ộ}{{\d \OCIRCUMFLEX }}1
		{Ổ}{{\h \OCIRCUMFLEX }}1
		{Ỗ}{{\~\OCIRCUMFLEX }}1
	%
	% Vần ơ
		{ơ}{{\ohorn }}1
		{ớ}{{\'\ohorn }}1
		{ờ}{{\`\ohorn }}1
		{ợ}{{\d \ohorn }}1
		{ở}{{\h \ohorn }}1
		{ỡ}{{\~\ohorn }}1
		%
		{Ơ}{{\OHORN }}1
		{Ớ}{{\'\OHORN }}1
		{Ờ}{{\`\OHORN }}1
		{Ợ}{{\d \OHORN }}1
		{Ở}{{\h \OHORN }}1
		{Ỡ}{{\~\OHORN }}1
	%
	% Vần u
		{ú}{{\'u}}1
		{ù}{{\`u}}1
		{ụ}{{\d u}}1
		{ủ}{{\h u}}1
		{ũ}{{\~u}}1
		%
		{Ú}{{\'U}}1
		{Ù}{{\`U}}1
		{Ụ}{{\d U}}1
		{Ủ}{{\h U}}1
		{Ũ}{{\~U}}1
	%
	% Vần ư
		{ư}{{\uhorn }}1
		{ứ}{{\'\uhorn }}1
		{ừ}{{\`\uhorn }}1
		{ự}{{\d \uhorn }}1
		{ử}{{\h \uhorn }}1
		{ữ}{{\~\uhorn }}1
		%
		{Ư}{{\UHORN }}1
		{Ứ}{{\'\UHORN }}1
		{Ừ}{{\`\UHORN }}1
		{Ự}{{\d \UHORN }}1
		{Ử}{{\h \UHORN }}1
		{Ữ}{{\~\UHORN }}1
	%
	% Vần y
		{ý}{{\'y}}1
		{ỳ}{{\`y}}1
		{ỵ}{{\d y}}1
		{ỷ}{{\h y}}1
		{ỹ}{{\~y}}1
		%
		{Ý}{{\'Y}}1
		{Ỳ}{{\`Y}}1
		{Ỵ}{{\d Y}}1
		{Ỷ}{{\h Y}}1
		{Ỹ}{{\~Y}}1
        %parentheses
        %
        {\{}{{\textcolor{parenthesesblue}{\{}}}{1}
        {\}}{{\textcolor{parenthesesblue}{\}}}}{1}
        {(}{{\textcolor{parenthesesblue}{(}}}{1}
        {)}{{\textcolor{parenthesesblue}{)}}}{1}
        {[}{{\textcolor{parenthesesblue}{[}}}{1}
        {]}{{\textcolor{parenthesesblue}{]}}}{1}
        % level 2
        {*\{}{{\textcolor{stringgreen}{\{}}}{1}
        {*\}}{{\textcolor{stringgreen}{\}}}}{1}
        {*(}{{\textcolor{stringgreen}{(}}}{1}
        {*)}{{\textcolor{stringgreen}{)}}}{1}
        {*[}{{\textcolor{stringgreen}{[}}}{1}
        {*]}{{\textcolor{stringgreen}{]}}}{1}
        % level 3
        {**\{}{{\textcolor{delimorange}{\{}}}{1}
        {**\}}{{\textcolor{delimorange}{\}}}}{1}
        {**(}{{\textcolor{delimorange}{(}}}{1}
        {**)}{{\textcolor{delimorange}{)}}}{1}
        {**[}{{\textcolor{delimorange}{[}}}{1}
        {**]}{{\textcolor{delimorange}{]}}}{1}
        %number
        {1}{{\textcolor{variableorange}{1}}}{1}
        {2}{{\textcolor{variableorange}{2}}}{1}
        {3}{{\textcolor{variableorange}{3}}}{1}
        {4}{{\textcolor{variableorange}{4}}}{1}
        {5}{{\textcolor{variableorange}{5}}}{1}
        {6}{{\textcolor{variableorange}{6}}}{1}
        {7}{{\textcolor{variableorange}{7}}}{1}
        {8}{{\textcolor{variableorange}{8}}}{1}
        {9}{{\textcolor{variableorange}{9}}}{1}
        {0}{{\textcolor{variableorange}{0}}}{1}
}


\begin{document}

% ---- Phần đầu: bìa, lời nói đầu, tóm tắt, mục lục ----
% Trang bìa. Thông tin môn học/đề tài lấy từ main.tex.
% Chỉ cần sửa danh sách giảng viên và sinh viên bên dưới.
\newgeometry{top=2cm,bottom=2cm,left=1.5cm,right=1.5cm,marginparwidth=2cm}
\pagenumbering{roman}   % bìa đánh số La Mã (không hiện) để mã neo trang không trùng với trang 1
\begin{titlepage}
    \linespread{1}\selectfont
    \begin{center}
        % Tên trường
        \leavevmode
        \normalfont
        {\LARGE \MakeUppercase{\head}}\\[0.3cm]
        {\Large \MakeUppercase{\university}}\\[0.3cm]
        {\large \MakeUppercase{\faculty}}
        \vskip 1cm
        % Logo trường
        \includegraphics[scale=0.2]{images/hcmus-logo.png}
        \vspace{1cm}

        \textbf{\huge \subjectname\ (\subjectcode)}\\[.3cm]
        \rule{\linewidth}{0.5mm}\\[.3cm]
        \textbf{\Large \reportname}\\[1.5cm]
        \textbf{\LARGE \reporttitle}\\[.3cm]
        \rule{\linewidth}{0.5mm}
        \\[.8cm]
        % Lớp và học kỳ (xoá dòng này nếu không cần)
        {\large Lớp: \classname \quad -- \quad \semester}
        \vfill

        % Giảng viên và sinh viên (mỗi người một dòng; thêm/bớt tuỳ ý)
        \begin{tabular}{r@{\hspace{1em}}l@{\hspace{1.5em}}l}
            \large \emph{Giảng viên hướng dẫn:} & \large Họ và tên giảng viên 1 & \\
            & \large Họ và tên giảng viên 2 & \\[1.2em]
            \large \emph{Sinh viên thực hiện:} & \large Họ và tên sinh viên 1 & \large (MSSV) \\
            & \large Họ và tên sinh viên 2 & \large (MSSV) \\
            & \large Họ và tên sinh viên 3 & \large (MSSV) \\
        \end{tabular}
        \vfill

        % Ngày tháng
        {\small \reportdate}
    \end{center}
\end{titlepage}
\restoregeometry
\pagenumbering{arabic}

% !TEX root = ../main.tex
\section*{Lời nói đầu}
Nội dung lời nói đầu và lời cảm ơn.

\medskip
\begin{flushright}
  \textbf{Nhóm thực hiện}
\end{flushright}

\clearpage
% !TEX root = ../main.tex
\section*{Tóm tắt}
Nội dung.

\clearpage
{\linespread{1}\selectfont\tableofcontents}
% Danh mục hình / bảng (tuỳ chọn): bỏ dấu % ở hai dòng dưới khi báo cáo có hình, bảng.
% {\linespread{1}\selectfont\listoffigures}
% {\linespread{1}\selectfont\listoftables}
\clearpage

% ---- Nội dung chính (đánh số mục) ----
% Muốn thêm phần: tạo file mới trong content/ rồi thêm một dòng \input tương ứng.
% !TEX root = ../main.tex
\section{Giới thiệu}
Nội dung.

% !TEX root = ../main.tex
\section{Cơ sở lý thuyết}
Nội dung.

% !TEX root = ../main.tex
\section{Nội dung thực hiện}
Nội dung.

% !TEX root = ../main.tex
\section{Kết quả và đánh giá}
Nội dung.

% !TEX root = ../main.tex
\section{Kết luận và hướng phát triển}
Nội dung.


% ---- Phần cuối: phụ lục, tài liệu tham khảo (mỗi phần một trang) ----
\clearpage
% !TEX root = ../main.tex
\section*{Phụ lục}
\addcontentsline{toc}{section}{Phụ lục}

\subsection*{Bảng phân công công việc}
\begin{table}[ht]
    \centering
    \begin{tabularx}{\textwidth}{|c|p{4.2cm}|X|c|}
        \hline
        \textbf{STT} & \textbf{Họ và tên (MSSV)} & \textbf{Công việc đảm nhận} & \textbf{Hoàn thành} \\
        \hline
        1 & Họ và tên (MSSV) & Mô tả công việc & 100\% \\
        \hline
        2 & Họ và tên (MSSV) & Mô tả công việc & 100\% \\
        \hline
    \end{tabularx}
\end{table}

\subsection*{Tài nguyên đính kèm}
Mã nguồn, dữ liệu, video demo: \url{https://...}

\clearpage
% !TEX root = ../main.tex
% Tài liệu tham khảo lấy từ references.bib. Trích dẫn trong văn bản bằng \cite{khoa}.
% Dòng \nocite dưới đây chỉ để template không báo lỗi khi chưa có \cite nào;
% xoá nó khi báo cáo đã trích dẫn tài liệu thật.
\nocite{example}
\addcontentsline{toc}{section}{Tài liệu tham khảo}
\bibliographystyle{unsrt}
\bibliography{references}


\end{document}
